\clearpage
\section{Conclusion} \label{sec:conclusion}
  This thesis examined the resilience of contemporary Raspberry Pi platforms against 
  disturbance-based memory faults. In particular, it investigated whether the Raspberry 
  Pi 4 and Raspberry Pi 5 exhibit susceptibility to Rowhammer-induced bit flips when 
  exposed to controlled hammering access patterns. The objective was to determine the 
  extent to which these modern embedded systems remain vulnerable to software-triggered 
  \gls{dram} disturbance errors.
  \\
  % Central Findings and Research Question
  The conducted experiments yielded several central findings. With regard to cache 
  bypass mechanisms, two out of four investigated approaches proved to be feasible. 
  Cache bypass could be achieved either through eviction strategies or by flushing 
  the cache using maintenance instructions. These methods enabled uncached memory 
  accesses required for effective hammering attempts.
  \\
  Concerning \gls{dram} mapping, the DRAMA tool~\cite{toolDrama} was successfully ported to 
  the ARM architecture. The adapted implementation provided accurate \gls{dram} address 
  mappings for the respective test platforms.
  \\
  In terms of outrunning the refresh mechanism, the TRRespass tool~\cite{toolTRRespass} 
  was likewise ported successfully to ARM systems. Although the tool operated as intended, 
  it does not identify any effective access patterns for the test devices.
  \\
  % Answer of the research question to what extent it was concluded -> Limitations
  Based on the experimental evaluation, it was concluded that modern Raspberry Pi 
  models, including Raspberry Pi 4 and Raspberry Pi 5 with RAM configurations ranging 
  from 4~GB to 16~GB, did not exhibit susceptibility to Rowhammer under the tested 
  conditions. Nevertheless, certain limitations prevent a definitive conclusion regarding 
  absolute immunity. In particular, the potential applicability of advanced 
  techniques such as Blacksmith~\cite{toolBlacksmith} and the influence of the \gls{dram} page 
  policy restrict the scope of the conclusion.
  \\
  % Comparison with Expectations and Prior Literature
  The absence of observable bit flips contrasts with prior expectations. Earlier studies 
  indicated that previous models, such as the Raspberry Pi 3, were susceptible to Rowhammer 
  attacks~\cite{paperRHonRi3B+}, and \gls{gpu}-based hammering~\cite{paperOpenGL} on the Raspberry 
  Pi 4. The lack of observable bit flips in the current study is likely due to active 
  \gls{trr} mechanisms, which mitigate disturbance effects and prevent reliable fault induction.
  \\
  % Implications for Security
  The findings narrow the immediate Rowhammer threat landscape for modern Raspberry Pi 
  platforms. As of the time of this study, experimental results indicate that these 
  devices are effectively resilient to conventional software-induced Rowhammer attacks.
  \\
  % Contribution and Scientific Value
  This work provides the first systematic evaluation of Rowhammer susceptibility on the 
  Raspberry Pi 5 across multiple RAM configurations, including 4~GB, 8~GB, and 16~GB variants. 
  It empirically demonstrates that two practical cache-bypass methods are implementable on 
  the Raspberry Pi 4 and Raspberry Pi 5. Furthermore, despite achieving successful cache 
  bypass and accurate \gls{dram} mapping, no effective hammering patterns could be induced.
  \\
  The results offer indirect evidence for the effectiveness of \gls{trr} as a mitigation 
  mechanism on ARM-based single-board computers. The study extends prior research that 
  focused primarily on the Raspberry Pi 3 and partially on the Raspberry Pi 4. It 
  provides clear experimental evidence that, under the applied conditions, no 
  reproducible bitflips occur.
  \\
  In addition, the a \gls{dram} Address Function reverse engineer approach using foundation of 
  \cite{paperDRAMA} was successfully ported to ARM systems 
  and specifically adapted for the Raspberry Pi platform. Similarly, a fully ported and 
  published version of the TRRespass tool~\cite{toolTRRespass} for 
  ARM systems was developed, thereby enabling further evaluation and experimentation.
  \\
  % Future Work
  Several open questions and potential extensions arise from this work.
  \\
  First, advanced investigations into \gls{trr} bypass techniques should be conducted. This 
  includes integrating and evaluating Blacksmith-style frequency-based hammering. 
  Furthermore, vendor-specific \gls{trr} behavior on \gls{lpddr} memory should be analyzed to 
  determine mitigation thresholds and activation logic. The development of adaptive 
  hammering techniques aimed at bypassing \gls{trr} and generating effective disturbance 
  patterns represents another important research direction.
  \\
  Second, further analysis of the page policy and memory controller behavior is 
  required. The active \gls{dram} page policy could be characterized using ~\cite{paperNethammer}.
  Additionally, the impact of the page policy on row activation frequency, 
  row buffer behavior, and overall disturbance feasibility should be evaluated. 
  Finally, methods to influence or manipulate page policy behavior should be 
  investigated in order to assess worst-case scenarios.
